\section{Conclusion}
\label{sec:conclusion}

This study compared the performance between two types of rendering in web applications, using Google's Core Web Vitals metrics as a quality assessment approach. For the analysis, both CSR and SSR rendering strategies were evaluated. A photo gallery application was developed for each rendering approach, and aspects such as performance, user experience, and memory usage were compared.

For data collection, a Node.js script was created using the Lighthouse performance measurement tool. The Kruskal-Wallis statistical test was applied to compare the results and verify the presence of statistically significant differences. The tests indicated relevant differences. One hypothesis raised is that CSR rendering provides better performance in metrics related to the loading of visible content, such as FCP and SI, while SSR rendering presents better results in metrics related to interactivity and loading of the largest page content, such as TTI and LCP. The results suggest that the choice between CSR and SSR should consider the specific objectives of the application and its usage environment.

\novo{Esta versão estendida aprofunda substancialmente a análise original ao introduzir um design experimental fatorial que incorpora duas técnicas de otimização de desempenho — minificação de JSON e compressão HTTP via gzip — como variáveis independentes. O experimento ampliado produziu 300 observações distribuídas em 30 condições de tratamento, possibilitando uma análise sistemática de como as técnicas de otimização interagem com a estratégia de renderização e o tamanho do payload.

Os resultados produzem três achados principais. Primeiro, a compressão gzip é a técnica de otimização de maior impacto dentre as avaliadas, reduzindo o LCP mediano do CSR em até 93,3\% para payloads Extra-Large (de 127,39~s para 8,58~s) e do SSR em até 95,0\% (de 29,35~s para 1,46~s). Segundo, a minificação de JSON no CSR proporciona um benefício significativo, porém comparativamente menor — reduzindo o LCP mediano em aproximadamente 47--52\% em payloads de médio a grande porte quando aplicada sem gzip; sua contribuição marginal quando combinada com o gzip é limitada (aproximadamente 19\% de melhoria adicional no tamanho Extra-Large). Terceiro, o SSR + Gzip supera consistentemente todas as configurações CSR em LCP e TTI, com a vantagem crescendo com o tamanho do payload: o SSR + Gzip é 2,3$\times$ mais rápido que o CSR + Minificação + Gzip no Extra-Small e 4,8$\times$ mais rápido no Extra-Large. Essa divergência é confirmada por tamanhos de efeito (Delta de Cliff) de $\delta = 1{,}0$ em todos os tamanhos de conjunto de dados para essa comparação, indicando separação prática completa entre as distribuições~\cite{Cliff1993,Romano2006}.

Um achado arquiteturalmente relevante diz respeito à invariância do FCP no CSR: a habilitação de gzip ou minificação no payload JSON não afeta o FCP, que permaneceu constante em aproximadamente 1,58~s independentemente do nível de otimização. Isso demonstra que o FCP no CSR é determinado pelo tempo de carregamento do script inicial e não é uma métrica adequada para avaliar otimizações de transferência de dados. LCP e TTI são as métricas apropriadas para esse fim, confirmando a orientação de seleção de métricas de~\citet{verma:2024}.

Do ponto de vista prático, esses achados fornecem orientações baseadas em evidências para arquitetos de aplicações web. A habilitação da compressão gzip é fortemente recomendada para qualquer aplicação web que utilize SSR ou CSR, dado que proporciona a maior melhoria de desempenho a um custo de implementação negligenciável na camada de aplicação~\cite{sikder:2016}. A minificação de JSON é recomendada como técnica complementar para aplicações CSR, com a compreensão de que seu benefício é maior quando o gzip não está disponível. Para aplicações onde LCP e TTI são as principais preocupações de qualidade — como plataformas de conteúdo servindo conjuntos de dados de médio a grande porte~\cite{gangishetti:2026} —, o SSR + Gzip proporciona uma vantagem clara e de grande efeito sobre o CSR independentemente do nível de otimização aplicado. Para aplicações onde o feedback visual inicial (FCP) é crítico e os payloads são pequenos, o CSR mantém desempenho competitivo e obtém pontuações Lighthouse compostas superiores.}

As future work, we suggest \novo{avaliar o efeito simétrico da minificação de HTML sobre o desempenho do SSR — fator não investigado neste estudo —, de modo a possibilitar uma comparação completa dos efeitos de minificação em ambas as estratégias de renderização. Adicionalmente, recomenda-se} analyzing other rendering approaches in web applications, such as Incremental Static Regeneration (ISR) and isomorphic architecture, would be of interest. It is also recommended to perform tests on mobile devices and under varying bandwidth conditions, given the relevance of such access scenarios to ensure an optimized and consistent experience across different platforms and contexts. \novo{Em particular, testes com limitação de rede simulada (por exemplo, condições 3G e 4G) esclareceriam se a vantagem do gzip é ainda mais amplificada em ambientes com restrição de largura de banda. Investigar o custo de CPU do gzip no servidor sob alta concorrência completaria a análise de compensação iniciada neste estudo. Além disso, replicar este estudo com frameworks SSR e CSR de produção (e.g., Next.js, Nuxt, SvelteKit, React) em múltiplos motores de navegador endereçaria a limitação de validade de construto identificada na Seção~\ref{sec:threats_to_validity}, conforme sugerido por~\citet{ollila:2021}.} Additionally, future research should include a broader range of application encompassing more interactive and diverse scenarios such as e-commerce platforms, dashboards, and complex single-page applications, in order to strengthen the generalizability of the findings.


\section*{ARTIFACT AVAILABILITY}
The \textit{scripts} and \textit{dataset} are available at:
\url{https://github.com/CleitonSilvaT/cs-vs-ssr-sbqs-2025}
